\section*{Bab \ref{ch:b3:crystalline-solids} --- Zat Padat Kristal}

\begin{solution}{exo:b3:crystalline-solids:1}
(a) sc: $8\times\tfrac18 = 1$; bcc: $8\times\tfrac18 + 1 = 2$;
fcc: $8\times\tfrac18 + 6\times\tfrac12 = 4$. (b) Diagonalnya
$a\sqrt3 = 4r$: $r = a\sqrt3/4$. (c) Diagonal mukanya $a\sqrt2 = 4r$:
$r = a\sqrt2/4$. (d) $0.68$ dan $0.74$: \omterm{rem:b3:crystalline-solids:bonds}{ikatan logam} yang buta arah
menginginkan kemasan yang maksimum, sehingga muncullah struktur fcc
(atau struktur heksagonal yang sama rapatnya) pada tembaga, aluminium,
perak dan emas.
\end{solution}

\begin{solution}{exo:b3:crystalline-solids:2}
(a) $a/\sqrt2 = \qty{2.55}{\angstrom}$. (b) 12 --- yakni
bilangan koordinasi kemasan rapatnya. (c) \qty{8960}{kg/m^3}
(\cref{ex:b3:crystalline-solids:density}). (d) $n = 4/a^3 =
\qty{8.5e28}{m^{-3}}$, sehingga satu mikrometer kubik menahan
$8.5\times10^{10}$ atom --- karena itulah metalurgi cip adalah
statistika, bukan pertukangan kayu.
\end{solution}

\begin{solution}{exo:b3:crystalline-solids:3}
(a) $\sin\theta_1 = \lambda/2d = 0.273$: $\theta_1 =
15.8^\circ$. (b) $n \le 2d/\lambda = 3.66$: tiga orde. (c)
$\sin\theta \le 1$ memaksa $n\lambda \le 2d$; panjang gelombang cahaya
tampak $\lambda \sim \qty{5000}{\angstrom}$ seribu kali terlalu besar
--- tak ada jarak bidang \omterm{def:b3:crystalline-solids:lattice}{kristal} yang dapat mendifraksikannya. (d) $\sin\theta$
turun 1\,\%: $\theta_1 \to 15.7^\circ$. Sebuah difraktometer adalah
termometer yang halus --- dan geseran inilah cara pemuaian termal
diukur pada skala atom.
\end{solution}

\begin{solution}{exo:b3:crystalline-solids:4}
(a) $d_{100} = a > d_{110} = a/\sqrt2 > d_{111} = a/\sqrt3$. (b)
$d$ terbesar, sudut terkecil: (100). (c) Keluarga oktahedral (111)
--- dalam struktur intan lembarannya berpasangan menjadi lapis ganda
yang berikatan kuat dengan celah lebar yang berikatan jarang di antaranya:
yakni bidang sang pembelah. (d) Tiap butirannya memantul pada $\theta$ yang sama
tetapi pada azimut yang acak: sinar pantulnya mengipas menjadi kerucut
bersetengah sudut $2\theta$, lalu memotong pendeteksinya sebagai cincin.
\end{solution}

\begin{solution}{exo:b3:crystalline-solids:5}
(a) Perpaduannya memerlukan selisih lintasan berorde $\lambda$;
jaraknya berorde ångström, sehingga $\lambda$ harus demikian pula. (b) $E =
hc/\lambda \approx \qty{8}{keV}$. (c) $p = h/\lambda$, $V =
p^2/2m_{\text{e}}e \approx \qty{60}{kV}$ (kenisbiannya memangkas beberapa
persen) --- yakni tegangan kerja sebuah mikroskop elektron. (d)
$\lambda = h/\sqrt{3m_{\text{n}}k_{\text{B}}T} \approx
\qty{1.5}{\angstrom}$: neutron bersuhu ruang lahir dalam keadaan siap
mendifraksi. Mereka menghambur dari intinya, bukan dari awan elektronnya
--- sehingga mereka melihat hidrogen dengan jelas dan membawa momen magnet
yang memetakan keteraturan magnetnya, dan keduanya nyaris tak terlihat sinar-X.
\end{solution}

\begin{solution}{exo:b3:crystalline-solids:6}
(a) Empat pasang tiap selnya yang bervolume $a^3$. (b) $N_{\text{A}} =
4M/\rho a^3 = 4\times0.05844/(2165\times\qty{1.794e-28}{})
\approx \qty{6.02e23}{mol^{-1}}$. (c) $N_{\text{A}} \propto
a^{-3}$: ralat 0.1\,\% pada $a$ menjadi 0.3\,\% pada $N_{\text{A}}$. (d)
Metodenya mencacah atom dengan menganggap tiap selnya penuh dan
serupa: kekosongan, pengotor dan batas mosaiknya semuanya
salah cacah --- sehingga diperlukan bola silikon yang sempurna secara
fanatik untuk pendefinisian ulang kilogram itu.
\end{solution}

\begin{solution}{exo:b3:crystalline-solids:7}
(a) Tiap ionnya melihat $2\sum_{n\ge1}(-1)^{n+1}/n = 2\ln2$ dalam satuan
$e^2/4\pi\varepsilon_0r_0$ (dengan faktor 2: kedua sisinya), dan bersifat menarik. (b)
$e^2/4\pi\varepsilon_0r_0 = \qty{5.11}{eV}$; $\times1.7476 =
\qty{8.93}{eV}$. (c) Peminimumannya membunuh $1/9$ tarikannya:
$U = 8.93\times\tfrac89 = \qty{7.94}{eV}$. (d) Dalam satu persen
dari \qty{7.9}{eV}: yakni \omterm{def:b3:crystalline-solids:lattice}{kisi} muatan titik ditambah satu eksponen
kekakuan menjelaskan seluruh anggaran sebuah zat padat ionik --- mekanika
kuantumnya bersembunyi di dalam $r_0$ dan eksponennya.
\end{solution}

\begin{solution}{exo:b3:crystalline-solids:8}
(a) $r_0 = 2^{1/6}\times3.40 = \qty{3.82}{\angstrom}$, yakni 1.5\,\%
di atas nilai terukurnya \qty{3.76}{\angstrom} (tiap atomnya juga merasakan
kedua belas tetangganya, yang mengetatkan sumurnya). (b)
$T_{\text{m}} \sim \epsilon/2k_{\text{B}} \approx \qty{60}{K}$:
yakni skala yang tepat untuk \qty{84}{K}. (c) \omterm{thm:b3:harmonic-oscillator:spectrum}{Energi titik nol} helium
$\hbar\omega/2$ dalam sumur yang begitu dangkal menyaingi sumurnya sendiri:
\omterm{def:b3:crystalline-solids:lattice}{kristalnya} mengguncang dirinya sampai berantakan, dan helium tetap cair pada $T = 0$
kecuali bila diperas. (d) Kedalaman $\sim\qty{0.01}{eV}$ berbanding kovalen
$\sim\qty{3.6}{eV}$ tiap ikatannya: tiga ratus kali lipat dalam energi adalah
seluruh jarak dari embun beku sampai intan.
\end{solution}

\begin{solution}{exo:b3:crystalline-solids:9}
(a) $ka \ll 1$: $\omega \approx 2\sqrt{K/m}\,(ka/2) =
a\sqrt{K/m}\,k$. (b) Rentangkanlah sebuah sel sebesar $\delta$: tegasannya $\sim
K\delta/a^2$, regangannya $\delta/a$, sehingga $E \sim K/a$, yakni $K \sim
Ea$. (c) $K \approx \qty{31}{N/m}$, $v = a\sqrt{K/m} \approx
\qty{4400}{m/s}$ --- yakni nilai terukurnya \qty{4000}{m/s} dalam 10\,\%,
dari pegas yang diterka dari \omterm{thm:b3:continuum-elasticity:hooke}{modulus Young}. (d) $v = a\sqrt{K/m}$:
ikatan yang kaku menaikkannya, atom yang ringan menaikkannya --- intan
memaksimumkan keduanya, sehingga bunyinya \qty{18000}{m/s} dan $\Theta_{\text{D}} \approx
\qty{2200}{K}$.
\end{solution}

\begin{solution}{exo:b3:crystalline-solids:10}
(a) $\dd\omega/\dd k \propto \cos(ka/2) = 0$ pada $k = \pi/a$. (b)
$u_n \propto (-1)^n$: tiap atomnya berlawanan fase dengan kedua
tetangganya --- yakni gelombang tegak, dengan energi yang berkecipak di tempat. (c)
Dengan $\lambda = 2a$, gelombang yang terhambur ke belakang oleh atom
berturutannya berselisih lintasan tepat satu panjang gelombang: yakni syarat Bragg
sepanjang rantainya sendiri. \omterm{def:b3:crystalline-solids:lattice}{Kisinya} memantulkan getarannya sendiri,
gelombang maju dan mundurnya terkunci menjadi pola tegak, dan
gelombang berjalannya tak dapat meneruskan perjalanannya. (d) $\omega_{\text{max}} = 2\sqrt{K/m} \approx
\qty{3.4e13}{rad/s}$, yakni $\sim\qty{5}{THz}$: yaitu inframerah jauh
--- getaran \omterm{def:b3:crystalline-solids:lattice}{kisi} dan cahaya inframerah bertemu dalam oktaf yang sama,
dan itulah kenyataan di balik latihan berikutnya.
\end{solution}

\begin{solution}{exo:b3:crystalline-solids:11}
(a) Dalam ragam optiknya, subkisi $+$ dan $-$ bergerak ke arah yang
berlawanan: yakni dwikutub listrik yang berayun, persis apa yang dicengkeram
gelombang cahaya. (b) $\omega \approx \sqrt{2K/\mu}$ dengan $\mu =
\qty{2.3e-26}{kg}$: $\omega \approx \qty{4.7e13}{rad/s}$,
$\lambda = 2\pi c/\omega \approx \qty{40}{\micro m}$ --- yakni inframerah
jauh (pita reststrahlen garam yang terukur duduk pada
$\qty{61}{\micro m}$: oktaf yang tepat dari model dua pegas). (c)
Pada pita itu talunan \omterm{def:b3:crystalline-solids:lattice}{kristalnya} sendiri menyerap lalu memantulkan kembali
gelombangnya: garam yang tembus berubah menjadi cermin yang kedap. (d) Satu atom tiap
selnya berarti tak ada subkisi kedua untuk saling berdetak: yakni satu cabang tunggal,
tanpa celah.
\end{solution}

\begin{solution}{exo:b3:crystalline-solids:12}
(a) $N$ atom $\times$ 3 arah simpangan = $3N$
osilator, sehingga spektrum lelurusnya dipenggal begitu ia telah
mencacah $3N$ ragam: itulah yang mendefinisikan $k_{\text{D}}$. (b)
$k_{\text{D}} = (6\pi^2n)^{1/3} = \qty{1.7e10}{m^{-1}}$,
$\Theta_{\text{D}} = \hbar vk_{\text{D}}/k_{\text{B}} \approx
\qty{350}{K}$ --- yakni nilai kalorimetrinya \qty{343}{K}, dari sebuah laju
bunyi dan sebuah kerapatan. (c) Garis lurusnya melewatkan pendataran
pada tepi zonanya: Debye mencacah berlebih frekuensi tingginya, sehingga
kecocokannya tegang pada suhu pertengahan; hukum $T^3$-nya (yakni hanya
gelombang panjang) tetap aman. (d) Keduanya mengatakan bahwa spektrumnya berakhir ketika panjang
gelombangnya mencapai jarak atomnya --- yakni satu ragam tiap atomnya, yang
berbusana entah sebagai vektor gelombang pemenggal atau sebagai suhu pemenggal.
\end{solution}

\begin{solution}{pb:b3:crystalline-solids:1}
\textbf{1.} $2d\sin\theta = n\lambda$. Ekawarna: satu
$\lambda$ membuat tiap sudut cincinnya memetakan ke satu jarak. Serbuk: di antara
butiran acaknya, selalu ada yang berarah tepat untuk memantul --- tiap
keluarganya berbicara sekaligus.
\textbf{2.} $\theta = 19.05^\circ, 22.15^\circ, 32.2^\circ,
38.75^\circ$; $\sin^2\theta = 0.107, 0.142, 0.284, 0.392$.
\textbf{3.} Sisipkanlah $d = a/\sqrt{h^2+k^2+l^2}$ ke dalam Bragg
($n = 1$): $\sin^2\theta = (\lambda^2/4a^2)(h^2+k^2+l^2)$.
\textbf{4.} Nisbahnya $1 : 1.33 : 2.67 : 3.68$; dikalikan 3:
$3.0 : 4.0 : 8.0 : 11.0$.
\textbf{5.} $(111)$ yang semuanya ganjil dan $(200)$, $(220)$,
$(311)$ yang semuanya genap memberikan tepat 3, 4, 8, 11; angka 1, 2, 5, 6, 7
yang absen menyingkirkan kubus sederhana, yang polanya akan menunjukkannya.
\textbf{6.} Jumlah genap yang diizinkan bcc memberikan nisbah
$1:2:3:4$ --- yakni irama yang berbeda. Nisbahnya selamat dari ketidaktahuan
akan $a$, dari hanyutan panjang gelombang, dan dari ketaksejajaran cuplikannya:
keganjilan adalah geometri, bukan kalibrasi.
\textbf{7.} Sedikit bidang membuat maksimum perpaduan yang tumpul, seperti
\omterm{def:b3:crystalline-solids:lattice}{kisi} berlima celah: lebar cincinnya $\sim$ kebalikan ukuran kristalitnya
(yakni hubungan Scherrer) --- cincin yang tajam mengesahkan butiran yang
tumbuh baik.
\textbf{8.} $a = \lambda\sqrt{h^2+k^2+l^2}/2\sin\theta = 4.088,
4.086, 4.089, 4.082\,\unit{\angstrom}$.
\textbf{9.} $a \approx \qty{4.09}{\angstrom}$.
\textbf{10.} $\rho = 4M/N_{\text{A}}a^3$.
\textbf{11.} $4/N_{\text{A}}a^3 = \qty{9.74e4}{mol/m^3}$:
perak $\to \qty{10.5}{g/cm^3}$ (cocok dengan nilai buku pegangan
perak), emas $\to 19.2$ (cocok dengan emas!), platina $\to 19.0$
(bertentangan dengan platina yang 21.4 --- $a$ sejatinya adalah
\qty{3.92}{\angstrom}). Platina tersingkir; perak dan emas keduanya
selamat, karena \omterm{def:b3:crystalline-solids:lattice}{tetapan kisinya} nyaris berimpit persis.
\textbf{12.} Pengukuran yang berdasar massa: yakni kerapatan (atau sekadar massa
molarnya). Perak dan emas berbagi $a$ sampai 0.2\,\%, tetapi massa atomnya
berbeda sebesar faktor 1.8 --- tak ada keberimpitan \omterm{def:b3:crystalline-solids:lattice}{kisi} yang dapat
menjembatani hal itu.
\textbf{13.} Archimedes pada lantakannya: di dekat \qty{10.5}{g/cm^3}
$\to$ perak, yang taat asas dengan karatnya (perak menghitam;
emas akan muncul berkilau). Massa sel mikroskopis dan
penimbangan makroskopisnya sepakat --- yakni
$\rho = 4M/N_{\text{A}}a^3$ yang sama yang dibaca dari kedua arahnya.
\textbf{14.} Atomnya bergemeretak di sekitar \emph{kedudukan rata-rata yang
tak bergeser}: \omterm{def:b3:crystalline-solids:lattice}{kisi} rata-ratanya, yang menetapkan sudut cincinnya, tetap utuh;
simpangan acaknya hanya membagi ulang keamatannya.
\textbf{15.} Keamatannya turun mulus bersama $T$ (yakni faktor
Debye--Waller bergaya $\exp(-q^2\langle x^2\rangle)$), dan cahaya yang
hilang muncul kembali sebagai latar baur di antara cincinnya.
\textbf{16.} $K \approx Ea = \qty{8.3e10}{}\times
\qty{2.89e-10}{} \approx \qty{24}{N/m}$.
\textbf{17.} $v = a\sqrt{K/m} \approx \qty{3300}{m/s}$ ---
yakni nilai terukurnya \qty{2700}{m/s} sampai 20\,\%.
\textbf{18.} $\langle x^2\rangle = 3k_{\text{B}}T/K$: akar rata-rata
kuadratnya $\approx\qty{0.23}{\angstrom}$, yakni kira-kira 8\,\% panjang
ikatannya sudah pada suhu ruang.
\textbf{19.} Tanda 10\,\%-nya terekstrapolasi ke $\sim\qty{500}{K}$
berbanding nilai sebenarnya \qty{1235}{K}: ordenya tepat, meleset kira-kira
dua kali lipat --- $K$ pegas tunggal kami terlalu lunak, dan Lindemann adalah
kaidah penskalaan, bukan hukum. Pelajarannya tetap: peleburan adalah ketika
gemeretak termalnya menyaingi \omterm{def:b3:crystalline-solids:lattice}{kisinya} sendiri.
\textbf{20.} $d_{110} = a/\sqrt2 = \qty{2.04}{\angstrom}$:
$\sin\theta = 0.378$, $2\theta \approx 44.5^\circ$.
\textbf{21.} Kekuatan gandengannya: elektron (yang bermuatan) menghambur dalam
nanometer --- yakni permukaan dan kerajang; sinar-X dalam mikrometer ---
yakni serbuk limbaknya; neutron dalam sentimeter --- yakni seluruh onderdil mesin.
\textbf{22.} Cincinnya runtuh menjadi satu atau dua lingkaran cahaya yang lebar:
cairannya menyimpan keteraturan jangkauan pendek (yakni jarak tetangga yang
disukainya) tetapi tanpa pencatatan jangkauan panjang --- tak ada lagi yang
berkala untuk berpadu secara tajam.
\textbf{23.} DNA --- yakni Foto 51. Difraksi membaca jarak
pengulangan jauh di bawah jangkauan mikroskop cahaya mana pun; pola X-nya
mengkhianati sebuah heliks, dan kaburan \qty{3.4}{\angstrom} itu adalah
anak tangganya yang bersusun.
\textbf{24.} Bahwa strukturnya tertata secara pasti
(kuasiberkala --- yakni tertata tanpa berulang): bintik tajam memerlukan
keserempakan fase jangkauan panjang, bukan kekalaan --- yakni penemuan
yang melebarkan definisi kristalografi sendiri tentang sebuah \omterm{def:b3:crystalline-solids:lattice}{kristal}.
\textbf{25.} Nisbahnya $3:4:8:11 \to$ fcc; $a = \qty{4.09}{\angstrom}$
$+$ kerapatannya $\to$ perak, bukan emas atau platina; dan karena pemanasan
memudarkan cincinnya dan peleburan menghapusnya, sinar-X-kanlah lantakannya
dalam keadaan dingin --- yakni perak museumnya, yang disahkan oleh perpaduan.
\end{solution}
